% AUDIT DE VERACITE 2026-09-20 -- corrections des citations de rapport de jury.
% Bloc 4, encadre jury : la citation attribuee au « Rapport jury 2023 » (« justifient explicitement les
% expressions litterales avant les applications numeriques ») est introuvable dans les 8 rapports
% disponibles (2018-2025) : citation fabriquee. REMPLACEE par le verbatim authentique du
% Rapport_Jury_2024.md ligne 515, de portee equivalente (exigence de justification des resultats),
% et re-attribuee en consequence au rapport 2024.
% La premiere citation (« Rapport jury 2024 », confusion couple/force) a ete verifiee conforme au
% Rapport_Jury_2024.md ligne 1079 et laissee intacte.
\expandafter\ifx\csname CAPETINCLUDEMODE\endcsname\relax
\documentclass[11pt,a4paper]{article}
% preamble.tex — packages et commandes communes a toutes les fiches CAPET SII IE
\input{glyphtounicode}
\pdfgentounicode=1
\usepackage{cmap}
\usepackage[T1]{fontenc}
\usepackage[utf8]{inputenc}
\usepackage[french]{babel}
% babel[french] active a \begin{document} des raccourcis de ponctuation (:;!?)
% qui cassent la syntaxe to[...] de circuitikz (erreur "+ or - expected") ;
% on les desactive, ce qui ne change que l'espacement fin avant ces signes.
\AtBeginDocument{\shorthandoff{:;!?}}
\usepackage[a4paper, top=1.6cm, bottom=1.6cm, left=1.4cm, right=1.4cm]{geometry}
\setlength{\headheight}{14.5pt}
\usepackage{amsmath, amssymb, amsfonts}
\usepackage{siunitx}
\sisetup{locale=FR, output-decimal-marker={,}, inter-unit-product=\ensuremath{{}\cdot{}},
         group-minimum-digits=4}
% Unites non fournies par siunitx, declarees ici une fois pour toutes (fiches + exercices).
% Les gardes \ifdefined evitent tout conflit avec une redeclaration locale restee dans un
% fichier : ne pas retirer les gardes, ne pas redeclarer ces unites dans les fichiers.
\DeclareSIUnit{\jour}{j}
\ifdefined\var\else\DeclareSIUnit{\var}{var}\fi   % puissance reactive
\ifdefined\vad\else\DeclareSIUnit{\vad}{VAD}\fi   % puissance deformante
\ifdefined\tr\else\DeclareSIUnit{\tr}{tr}\fi      % tour (tr/min)
\usepackage{array, booktabs, multirow, tabularx}
\usepackage{xcolor}
\usepackage[most, breakable]{tcolorbox}
\usepackage{multicol}
\usepackage{enumitem}
\usepackage{mathtools}
\usepackage{lmodern}
% Refonte visuelle 2026 (direction "moderne coloré") : Lato pour le texte,
% maths sans-serif assorties via newtxsf (charge apres amsmath/amssymb).
\usepackage[default]{lato}
\usepackage{newtxsf}
\usepackage{microtype}
\usepackage{fancyhdr}
\usepackage{tikz}
\usetikzlibrary{positioning,arrows.meta}
\usepackage{circuitikz}
\usepackage{lastpage}

% ── Couleurs par domaine (palette modernisée 2026) ────────────────────
\definecolor{colElec}{RGB}{29, 78, 216}        % Électrotechnique — bleu vif
\definecolor{colElecLight}{RGB}{219, 234, 254} % fond bleu clair
\definecolor{colAuto}{RGB}{21, 128, 61}        % Automatique — vert
\definecolor{colMeca}{RGB}{194, 65, 12}        % Mécanique — orange
\definecolor{colTherm}{RGB}{185, 28, 28}       % Thermique — rouge
\definecolor{colQT}{RGB}{232, 244, 233}        % (hérité, plus utilisé par questiontype)
\definecolor{colJury}{RGB}{255, 248, 225}      % fond jury (crème)
\definecolor{colJuryFrame}{RGB}{202, 138, 4}   % bordure/bandeau jury (ambre)
\definecolor{colGrayLight}{RGB}{245, 245, 247} % fond gris neutre
\definecolor{colSeparator}{RGB}{180, 180, 180} % gris séparateur
\definecolor{colExo}{RGB}{8, 145, 178}          % Banque d'exercices — cyan (distinct des fiches)
% ── Encre et gris de la refonte ───────────────────────────────────────
\definecolor{colInk}{RGB}{23, 32, 42}          % texte courant (quasi-noir)
\definecolor{colSlate}{RGB}{83, 98, 118}       % gris ardoise (méta, en-têtes)
\definecolor{colRule}{RGB}{203, 213, 225}      % filets clairs
\definecolor{colJuryBg}{RGB}{255, 248, 225}    % fond crème jury
\definecolor{colPresence}{RGB}{83, 98, 118}    % bloc 5 (présence dans les sujets)
\AtBeginDocument{\color{colInk}}

% ── Style encadré exercice + corrigé (banque d'exercices) ─────────────
\tcbset{
  exoBox/.style={
    breakable,
    enhanced,
    colback=colExo!8!white,
    colframe=colExo,
    colbacktitle=colExo,
    coltitle=white,
    fonttitle=\bfseries\small,
    left=4pt, right=4pt, top=3pt, bottom=3pt,
    before skip=5pt, after skip=5pt,
  }
}

% ── Bascule enonce / corrige (banque d'exercices) ─────────────────────
% \CAPETCORRIGE est defini en ligne de commande (pdflatex "\def\CAPETCORRIGE{1}\input{...}")
% pour compiler la variante corriges d'un meme fichier source ; par defaut (non defini),
% la compilation produit la variante enonces. \CAPETSHOW{enonce}{corrige} choisit la branche.
%
% TROIS MODES (refonte 2026-09-20, manuel de revision) :
%   - aucun drapeau   -> imprime l'enonce            (variante enonces, INCHANGE)
%   - \CAPETCORRIGE   -> imprime le corrige          (variante corriges, INCHANGE)
%   - \CAPETMANUEL    -> imprime l'enonce ET met le corrige de cote pour l'annexe
%
% Le mode manuel stocke chaque corrige dans une macro (\csgdef), sans ecriture
% fichier ni re-tokenisation : le contenu (circuitikz, accents, %) traverse intact.
% L'annexe les rejoue ensuite dans l'ordre via \CAPETANNEXE.
\usepackage{etoolbox}

\newcounter{CAPETexo}          % numero continu d'exercice dans le manuel
\setcounter{CAPETexo}{0}
\newcounter{CAPETannexe}       % compteur de relecture en annexe

% Titre courant de la notion en cours (pose par \CAPETNOTION), sert d'intertitre annexe
\newcommand{\CAPETnotioncode}{}
\newcommand{\CAPETnotiontitre}{}

% \CAPETNOTION{code}{titre} : ouvre une notion dans le manuel (memorise le libelle,
% et marque le rang du premier exercice de la notion pour l'intertitre d'annexe).
\newcommand{\CAPETNOTION}[2]{%
  \renewcommand{\CAPETnotioncode}{#1}%
  \renewcommand{\CAPETnotiontitre}{#2}%
  \csgdef{CAPETnotionstart\the\numexpr\value{CAPETexo}+1\relax}{#1 -- #2}%
}

% ── \degree robuste ───────────────────────────────────────────────────
% Certains fichiers d'exercices utilisent \degree. Ni le preambule ni siunitx
% ne le definissent, et \providecommand local ne suffit pas : sous hyperref,
% \degree existe deja comme commande factice (\ERROR), donc \providecommand
% s'abstient et l'appel echoue a la compilation du manuel (erreur invisible
% en compilation isolee). On le (re)definit donc inconditionnellement ici,
% apres le chargement des packages, pour tout le corpus.
\AtBeginDocument{\providecommand{\degree}{}\renewcommand{\degree}{\ensuremath{^\circ}}}

% Libelle decoratif : les fichiers d'exercices utilisent aussi \CAPETSHOW dans
% leur \fancyhead et leur bloc de titre, pour afficher "(enonces)" ou "(corriges)".
% Ces appels-la ne sont PAS des exercices : en mode manuel ils ne doivent ni
% consommer un numero, ni alimenter l'annexe. On les reconnait a leur premier
% argument, strictement egal au libelle ci-dessous.
\newcommand{\CAPETLIBELLEENONCES}{(\'enonc\'es)}

\ifdefined\CAPETMANUEL
  % --- Mode manuel : enonce en place, corrige differe vers l'annexe ---
  \newcommand{\CAPETSHOW}[2]{%
    \def\CAPETargun{#1}%
    \ifx\CAPETargun\CAPETLIBELLEENONCES
      #1%            appel decoratif : rendu tel quel, hors numerotation
    \else
      \CAPETSHOWEXO{#1}{#2}%
    \fi
  }
  \newcommand{\CAPETSHOWEXO}[2]{%
    \stepcounter{CAPETexo}%
    \csgdef{CAPETcorrige\theCAPETexo}{#2}%
    \edef\CAPETtmpnum{\theCAPETexo}%
    #1%
    \par\vspace{2pt}%
    {\footnotesize\itshape\color{colExo}%
     \hyperref[capet:cor:\CAPETtmpnum]{Corrige de l'exercice \CAPETtmpnum\ $\rightarrow$ p.~\pageref{capet:cor:\CAPETtmpnum}}%
     \label{capet:eno:\CAPETtmpnum}}%
  }
  % \CAPETANNEXE : rejoue tous les corriges captures, dans l'ordre, avec intertitres
  \newcommand{\CAPETANNEXE}{%
    \setcounter{CAPETannexe}{0}%
    \loop\ifnum\value{CAPETannexe}<\value{CAPETexo}%
      \stepcounter{CAPETannexe}%
      \edef\CAPETcurnum{\theCAPETannexe}%
      \ifcsname CAPETnotionstart\CAPETcurnum\endcsname
        \par\medskip
        {\normalsize\bfseries\color{colExo}\csuse{CAPETnotionstart\CAPETcurnum}}%
        \par\nobreak\smallskip
      \fi
      \begin{tcolorbox}[exoBox, breakable,
        title={Corrige de l'exercice \CAPETcurnum\
               \hfill\normalfont\footnotesize
               \hyperref[capet:eno:\CAPETcurnum]{enonce p.~\pageref{capet:eno:\CAPETcurnum}}}]%
      \label{capet:cor:\CAPETcurnum}%
      \csuse{CAPETcorrige\CAPETcurnum}%
      \end{tcolorbox}%
    \repeat
  }
\else
  % --- Modes historiques, comportement strictement inchange ---
  \newcommand{\CAPETSHOW}[2]{\ifdefined\CAPETCORRIGE #2\else #1\fi}
  \newcommand{\CAPETANNEXE}{}
\fi

% ── Difficulté 1 à 3 (étoiles pleines, amssymb) ───────────────────────
\newcommand{\difficulte}[1]{%
  \ifcase#1\relax
  \or $\bigstar$\;niveau 1/3%
  \or $\bigstar\bigstar$\;niveau 2/3%
  \or $\bigstar\bigstar\bigstar$\;niveau 3/3%
  \fi
}

% ── Style encadré extrait de document technique (banque d'exercices) ──
\tcbset{
  docTech/.style={
    breakable,
    enhanced,
    colback=colGrayLight,
    colframe=colSeparator,
    fonttitle=\bfseries\small\itshape,
    coltitle=black,
    colbacktitle=colGrayLight,
    left=4pt, right=4pt, top=3pt, bottom=3pt,
    before skip=4pt, after skip=4pt,
  }
}

% ── Style tcolorbox de section principale ────────────────────────────
% Refonte 2026 : bandeau titre coloré arrondi + fond très pâle, sans cadre.
\tcbset{
  ficheSection/.style={
    breakable, enhanced jigsaw,
    frame hidden, boxrule=0pt, arc=1.8mm,
    colback=#1!4!white,
    colbacktitle=#1,
    coltitle=white,
    fonttitle=\sffamily\bfseries,
    left=7pt, right=7pt, top=5pt, bottom=5pt,
    toptitle=3pt, bottomtitle=3pt,
    before skip=10pt, after skip=8pt,
  }
}

% ── Style encadré question type ──────────────────────────────────────
% Carte blanche, bandeau sombre, liseré très clair.
\tcbset{
  questiontype/.style={
    breakable, enhanced jigsaw,
    colback=white,
    frame hidden, boxrule=0pt, arc=1.8mm,
    colbacktitle=colElec!78!black,
    coltitle=white,
    fonttitle=\sffamily\bfseries\small,
    left=7pt, right=7pt, top=4pt, bottom=4pt,
    toptitle=2.5pt, bottomtitle=2.5pt,
    before skip=8pt, after skip=8pt,
    borderline={0.6pt}{0pt}{colElec!12!white},
  }
}

% ── Style encadré jury ───────────────────────────────────────────────
% Bandeau ambre, fond crème, sans cadre.
\tcbset{
  jury/.style={
    breakable, enhanced jigsaw,
    colback=colJuryBg,
    frame hidden, boxrule=0pt, arc=1.8mm,
    colbacktitle=colJuryFrame,
    coltitle=white,
    fonttitle=\sffamily\bfseries,
    left=7pt, right=7pt, top=4pt, bottom=4pt,
    toptitle=3pt, bottomtitle=3pt,
    before skip=10pt, after skip=8pt,
  }
}

% ── En-tête / pied de page communs (refonte 2026) ─────────────────────
% #1 code fiche, #2 titre court, #3 couleur domaine
\newcommand{\ficheHeaders}[3]{%
  \fancyhead[L]{\small\sffamily{\bfseries\textcolor{#3}{#1}}\hspace{0.7em}\textcolor{colSlate}{#2}}%
  \fancyhead[R]{\small\sffamily\textcolor{colSlate}{CAPET SII IE\, --- \,Session 2026}}%
  \fancyfoot[C]{\small\sffamily\textcolor{colSlate}{\thepage\,/\,\pageref{LastPage}}}%
  \fancyhead[C]{}\fancyfoot[L]{}\fancyfoot[R]{}%
}

% ── Bloc en-tête de fiche (BLOC 0, refonte 2026) ──────────────────────
% #1 code  #2 titre  #3 couleur domaine  #4 domaine  #5 ligne présence  #6 sous-thèmes
\newcommand{\ficheEntete}[6]{%
  \begin{tcolorbox}[enhanced, breakable,
    frame hidden, boxrule=0pt, arc=2.5mm,
    interior style={left color=#3, right color=#3!75!black},
    left=9pt, right=9pt, top=7pt, bottom=7pt,
    before skip=0pt, after skip=6pt]
    {\color{white}\sffamily\bfseries\LARGE #1\quad #2}\\[4pt]
    {\color{#3!22!white}\sffamily\small \MakeUppercase{#4}\quad\textbar\quad #5}\\[1pt]
    {\color{#3!35!white}\sffamily\footnotesize Sous-thèmes : #6}
  \end{tcolorbox}%
}

% ── Mise en page ─────────────────────────────────────────────────────
\pagestyle{fancy}
\fancyhf{}
\renewcommand{\headrulewidth}{0.5pt}
\renewcommand{\footrulewidth}{0pt}
\AtBeginDocument{%
  \renewcommand{\headrule}{{\color{colRule}\hrule width\headwidth height 0.5pt}}%
}

\setlist[itemize]{leftmargin=1.2em, itemsep=1pt, topsep=2pt, parsep=0pt}
\setlist[enumerate]{leftmargin=1.2em, itemsep=3pt, topsep=2pt, parsep=0pt}

\setlength{\parskip}{2pt}
\setlength{\columnsep}{10pt}

\def\CAPETBEGINFILE{\begin{document}}
\def\CAPETENDFILE{\end{document}}
\else
\def\CAPETBEGINFILE{}
\def\CAPETENDFILE{}
\fi

\ficheHeaders{T10}{Mécanique : PFD, réducteurs}{colMeca}

\CAPETBEGINFILE

% -----------------------------------------------------------------------
%  BLOC EN-TETE
% -----------------------------------------------------------------------

\ficheEntete{T10}{Mécanique : cinématique, forces, moments, réducteurs}{colMeca}%
  {Mécanique}%
  {\textbf{Présence :} \textbf{13/13} sujets (2013-2026) - thème \emph{systématique}}%
  {PFD translation/rotation - couple et force - rapport de réduction - inertie ramenée au moteur - rendement chaîne - couple thermique équivalent}

% -----------------------------------------------------------------------
%  1. RAPPELS DE COURS
% -----------------------------------------------------------------------

\begin{tcolorbox}[
  ficheSection=colMeca,
  title={1. Rappels de cours},
]

% -- 1.1 Principe fondamental de la dynamique -------------------------
\textbf{1.1 Principe fondamental de la dynamique (PFD)}

\smallskip
\begin{center}
\renewcommand{\arraystretch}{1.45}
\begin{tabular}{lll}
\toprule
\textbf{Mouvement} & \textbf{Équation} & \textbf{Unité} \\
\midrule
Translation & $\sum F = m \cdot \gamma$ & N, kg, m\,s$^{-2}$ \\[4pt]
Rotation axe fixe & $\sum M_O = J_{eq} \cdot \dot{\Omega}$ & N\,m, kg\,m$^2$, rad\,s$^{-2}$ \\[4pt]
Accélération constante & $V(t) = \gamma \cdot t$ ; \; $x(t) = \frac{1}{2}\gamma t^2$ & m\,s$^{-1}$, m \\
\bottomrule
\end{tabular}
\end{center}

\medskip
\hrule\vspace{4pt}

% -- 1.2 Couple, puissance et reducteur -------------------------------
\textbf{1.2 Couple, puissance et réducteur}

\smallskip
\[
  P_{meca} = C \cdot \Omega = F \cdot V
  \qquad;\qquad
  C_e = F \cdot r \quad (r = \text{bras de levier})
\]
Rapport de réduction : $r = \Omega_{sortie}/\Omega_{moteur} < 1$ (en général).
\[
  R_{tot} = \frac{N_{moteur}}{N_{charge}}
  \qquad;\qquad
  \eta_{tot} = \eta_{r1} \cdot \eta_{r2} \cdots \quad \text{(rendements en série)}
\]
Couple moteur en fonction de l'effort de traction (conservation des puissances) :
\[
  C_m = \frac{C_e \cdot r}{\eta_{red}} = \frac{F \cdot r_{roue} \cdot r_{red}}{\eta_{red}}
\]

\smallskip
\textbf{Chaîne cinématique motoréducteur :}
\begin{center}
\begin{tikzpicture}[>=Stealth, font=\small, node distance=10mm]
  \node[draw, align=center, minimum width=22mm, minimum height=11mm] (M) {Moteur\\$J_{mot}$};
  \node[draw, align=center, minimum width=20mm, minimum height=11mm, right=30mm of M] (R) {Réducteur\\$r$, $\eta_{red}$};
  \node[draw, align=center, minimum width=22mm, minimum height=11mm, right=30mm of R] (C) {Charge\\$J_{charge}$};
  \draw[->] (M) -- node[above, font=\footnotesize]{$\Omega_{mot}$} (R);
  \draw[->] (R) -- node[above, font=\footnotesize]{$\Omega_{sortie}=r\Omega_{mot}$} (C);
\end{tikzpicture}
\end{center}

\medskip
\hrule\vspace{4pt}

% -- 1.3 Inertie ramenee au moteur ------------------------------------
\textbf{1.3 Inertie équivalente ramenée à l'arbre moteur}

\smallskip
\[
  J_{eq} = J_{mot} + r^2 \cdot J_{charge}
  \qquad (r = \Omega_{sortie}/\Omega_{moteur})
\]
Théorème de Huygens (transfert d'axe) :
\[
  J_{(O,\delta)} = J_{(G,\delta)} + m \cdot d^2
\]
Moment d'inertie d'un parallélépipède par rapport à un axe central :
\[
  J_{(G,x)} = \frac{m}{12}(b^2 + c^2)
\]

\medskip
\hrule\vspace{4pt}

% -- 1.4 Couple thermique equivalent ----------------------------------
\textbf{1.4 Couple thermique équivalent}

\smallskip
Sur un cycle de période $T$ avec des phases de couple $C_{mi}$ de durée $\Delta t_i$ :
\[
  C_{th} = \sqrt{\frac{\sum_{i} C_{mi}^2 \cdot \Delta t_i}{T}}
\]
Condition de dimensionnement : $C_{th} \leq C_{nominal,moteur}$.

\end{tcolorbox}

% -----------------------------------------------------------------------
%  2. FORMULAIRE SYNTHETIQUE
% -----------------------------------------------------------------------

\begin{tcolorbox}[
  ficheSection=colMeca,
  title={2. Formulaire synthétique},
]
\small
\begin{multicols}{2}
\renewcommand{\arraystretch}{1.45}
\begin{tabular}{@{}p{4.5cm}ll@{}}
\toprule
\textbf{Formule} & \textbf{Usage} & \textbf{U.} \\
\midrule
\multicolumn{3}{@{}l}{\footnotesize\textit{PFD}} \\
$\sum F = m\gamma$ & Translation & N \\
$\sum M = J_{eq}\dot\Omega$ & Rotation & N\,m \\
$V = \gamma t$ & Accél. cste & m/s \\
$x = \frac{1}{2}\gamma t^2$ & Distance & m \\
\midrule
\multicolumn{3}{@{}l}{\footnotesize\textit{Couple et puissance}} \\
$P = C\Omega = FV$ & Puiss. méca. & W \\
$C_e = F \cdot r_{roue}$ & Couple roue & N\,m \\
$C_m = C_e \cdot r/\eta$ & Couple moteur & N\,m \\
$R_{tot} = N_{mot}/N_{charge}$ & Rapport réd. & - \\
$\eta_{tot} = \prod \eta_i$ & Rendement & - \\
\bottomrule
\end{tabular}

\columnbreak

\renewcommand{\arraystretch}{1.45}
\begin{tabular}{@{}p{4.5cm}ll@{}}
\toprule
\textbf{Formule} & \textbf{Usage} & \textbf{U.} \\
\midrule
\multicolumn{3}{@{}l}{\footnotesize\textit{Inertie}} \\
$J_{eq} = J_{mot} + r^2 J_{charge}$ & Ramené moteur & kg\,m$^2$ \\
$J_{(O)} = J_{(G)} + md^2$ & Huygens & kg\,m$^2$ \\
$J_{(G)} = m(b^2+c^2)/12$ & Parallélépipède & kg\,m$^2$ \\
\midrule
\multicolumn{3}{@{}l}{\footnotesize\textit{Couple thermique}} \\
$C_{th} = \sqrt{\sum C_i^2\Delta t_i/T}$ & Échauffement & N\,m \\
$C_{th} \leq C_{nom}$ & Condition & - \\
\midrule
\multicolumn{3}{@{}l}{\footnotesize\textit{Conversion}} \\
$\Omega = 2\pi N/60$ & tr/min → rad/s & rad/s \\
$C = P/\Omega$ & Couple nominal & N\,m \\
\bottomrule
\end{tabular}
\end{multicols}
\normalsize
\end{tcolorbox}

% -----------------------------------------------------------------------
%  3. QUESTIONS TYPES
% -----------------------------------------------------------------------

\begin{tcolorbox}[
  ficheSection=colMeca,
  title={3. Questions types - extraites des sujets réels},
]

% -- QT1 -------------------------------------------------------------

\begin{tcolorbox}[questiontype,
  title={QT1 - Couple moteur par conservation des puissances (rame MF67) \hfill\normalfont\footnotesize E2 sujet 2019 (Q25-27)}
]
\textbf{Énoncé.}
Chaque moteur CC est relié à l'essieu par un réducteur $\rho = 0{,}17$, $\eta_{red} = 97\,\%$. Diamètre de roue $d_e = 0{,}86\,\mathrm{m}$. Efforts de traction : $F_{e1} = 14{,}9\,\mathrm{kN}$ et $F_{e2} = 9{,}2\,\mathrm{kN}$.
\begin{enumerate}[label=(\alph*)]
  \item Exprimer le couple $C_e$ à l'essieu en fonction de $F_e$.
  \item Par conservation des puissances, exprimer le couple moteur $C_m$.
  \item Calculer $C_{m1}$ et $C_{m2}$.
\end{enumerate}

\tcbline

\textbf{Corrigé.}

\textbf{Étape 1 (a).} $C_e = F_e \cdot \dfrac{d_e}{2} = F_e \times 0{,}43\,\mathrm{m}$\\
$C_{e1} = 14\,900 \times 0{,}43 = 6\,407\,\mathrm{N\!\cdot\!m}$ ;\quad $C_{e2} = 9\,200 \times 0{,}43 = 3\,956\,\mathrm{N\!\cdot\!m}$

\textbf{Étape 2 (b).} Conservation des puissances : $C_e \cdot \Omega_{essieu} = \eta_{red} \cdot C_m \cdot \Omega_{moteur}$ avec $\Omega_{essieu}/\Omega_{moteur} = \rho$ :
\[
  C_m = \frac{C_e \cdot \rho}{\eta_{red}}
\]

\textbf{Étape 3 (c).}
\[
  C_{m1} = \frac{6\,407 \times 0{,}17}{0{,}97} \approx \mathbf{1\,122\,\mathrm{N\!\cdot\!m}}
  \qquad
  C_{m2} = \frac{3\,956 \times 0{,}17}{0{,}97} \approx \mathbf{693\,\mathrm{N\!\cdot\!m}}
\]

\textbf{Conclusion.} Les couples moteurs déterminés avec le rendement de réducteur pris en compte sont directement comparables au couple nominal du moteur retenu, conformément au cahier des charges de traction du cas étudié.
\end{tcolorbox}

\vspace{4pt}

% -- QT2 -------------------------------------------------------------

\begin{tcolorbox}[questiontype,
  title={QT2 - Inertie équivalente ramenée au moteur (bras robot KUKA) \hfill\normalfont\footnotesize E2 appliquée 2025 (Q9-11)}
]
\textbf{Énoncé.}
Structure plateau : $m = \SI{181}{\kilogram}$, dimensions $3000 \times 2400 \times \SI{60}{\milli\metre}$. Le pivot est en $O$ (bord du plateau, à $d = \SI{1,5}{\metre}$ du centre de gravité $G$). Rapport de réduction $r = 0{,}00607$, $J_{mot} = \SI{0,003}{\kilogram\metre\squared}$. Montrer que $J_{eq} = \SI{0,0118}{\kilogram\metre\squared}$.

\tcbline

\textbf{Corrigé.}

\textbf{Étape 1 — Inertie au centre de gravité.}
\[
  J_{(G,z)} = \frac{m}{12}(a^2 + b^2) = \frac{181}{12}(3{,}0^2 + 2{,}4^2) = \frac{181}{12}(9 + 5{,}76) = \SI{222,5}{\kilogram\metre\squared}
\]

\textbf{Étape 2 — Transfert par Huygens.}
\[
  J_{(O,z)} = J_{(G,z)} + m\,d^2 = 222{,}5 + 181 \times 1{,}5^2 = 222{,}5 + 407{,}25 \approx \SI{629}{\kilogram\metre\squared}
\]

\textbf{Étape 3 — Ramener au moteur.}
\[
  J_{eq} = J_{mot} + r^2 \cdot J_{(O,z)} = 0{,}003 + (0{,}00607)^2 \times 629
  = 0{,}003 + 0{,}0232 \approx \SI{0,026}{\kilogram\metre\squared}
\]

\textbf{Conclusion. [corrigé reconstitué]} Avec les seules données reproduites ici (distance $d \approx \SI{1,5}{\metre}$ estimée d'après le DT), la démarche littérale $J_{eq} = J_{mot} + r^2 J_{(O,z)}$ — qui est ce que le jury évalue — aboutit à $\approx \SI{0,026}{\kilogram\metre\squared}$, alors que le sujet officiel annonce $J_{eq} = \SI{0,0118}{\kilogram\metre\squared}$ (obtenu avec le $J_{(O,z)}$ précis calculé en Q8 du sujet, non disponible ici). Retenir la méthode (parallélépipède $\to$ Huygens $\to$ facteur $r^2$) et la valeur officielle \SI{0,0118}{\kilogram\metre\squared} comme référence, en signalant l'écart plutôt qu'en forçant la correspondance.
\end{tcolorbox}

\vspace{4pt}

% -- QT3 -------------------------------------------------------------

\begin{tcolorbox}[questiontype,
  title={QT3 - Freinage Urbanloop - vérification de décélération \hfill\normalfont\footnotesize E2 appliquée 2024 (Q17)}
]
\textbf{Énoncé.}
La capsule ($M_C = 1000\,\mathrm{kg}$) doit décélérer à $|a_C| = 1{,}4\,\mathrm{m\!\cdot\!s^{-2}}$. Les freins fournissent : $C_{FAV} = 150\,\mathrm{N\!\cdot\!m}$ par roue avant ($\times 2$), $C_{FAR} = 60\,\mathrm{N\!\cdot\!m}$ par roue arrière ($\times 2$), rayon $R_{roue} = 0{,}20\,\mathrm{m}$. Vérifier que la décélération est atteinte.

\tcbline

\textbf{Corrigé.}

\textbf{Étape 1 — Forces de freinage.}
\[
  T_{FAV} = \frac{150}{0{,}20} \times 2 = 1500\,\mathrm{N}
  \qquad
  T_{FAR} = \frac{60}{0{,}20} \times 2 = 600\,\mathrm{N}
  \qquad
  T_F = 2100\,\mathrm{N}
\]

\textbf{Étape 2 — Décélération.}
\[
  |a_C| = \frac{T_F}{M_C} = \frac{2100}{1000} = 2{,}1\,\mathrm{m\!\cdot\!s^{-2}}
\]

\textbf{Conclusion.} $2{,}1 > 1{,}4\,\mathrm{m\!\cdot\!s^{-2}}$ : le cahier des charges de freinage est respecté, avec une marge de sécurité confortable.\checkmark
\end{tcolorbox}

\end{tcolorbox}

% -----------------------------------------------------------------------
%  4. ATTENTES DU JURY
% -----------------------------------------------------------------------

\begin{tcolorbox}[
  jury,
  title={4. Attentes du jury - Erreurs récurrentes et points de vigilance},
]

\begin{enumerate}[label=\textbf{\arabic*.}, leftmargin=*, itemsep=5pt]

  \item \textbf{Confondre force et couple.}
  Un couple s'exprime en N\,m (force $\times$ bras de levier). Une force s'exprime en N.
  Le jury 2024 le signale explicitement comme erreur fréquente.

  \item \textbf{Oublier le rendement.}
  Lors du passage arbre charge → arbre moteur, le rendement $\eta$ divise : $C_m = C_{charge} \cdot r / \eta$.
  Ne pas l'omettre même si la valeur semble proche de 1.

  \item \textbf{Sens du rapport de réduction.}
  $r = \Omega_{sortie}/\Omega_{moteur} < 1$ (en général). Ne pas l'inverser.

  \item \textbf{Inertie ramenée : facteur $r^2$ (et non $1/r^2$).}
  $J_{eq} = J_{mot} + r^2 \cdot J_{charge}$. Une grande réduction donne une faible contribution de la charge côté moteur.

  \item \textbf{$v = d/t$ invalide en phase d'accélération.}
  Cette relation ne s'applique que si la vitesse est constante. Utiliser $V(t) = \gamma t$ et $x(t) = \frac{1}{2}\gamma t^2$.

  \item \textbf{Énoncer les hypothèses.}
  Le jury valorise les candidats qui précisent les hypothèses (rendement idéal, glissement nul, inertie négligée) avant le calcul.

  \item \textbf{Toujours conclure par rapport au cahier des charges.}
  Comparer le résultat à la valeur imposée et conclure explicitement (valide / non valide).

\end{enumerate}

\medskip
\begin{tcolorbox}[
  colback=colJuryFrame!12!colJury,
  colframe=colJuryFrame,
  left=4pt, right=4pt, top=2pt, bottom=2pt,
  fontupper=\small\itshape
]
Rapport jury 2024 : «~Le jury regrette la confusion trop fréquente entre couple et force et rappelle que la formule $v = d/t$ n'est applicable que si la vitesse est constante. Il invite les candidats à vérifier l'homogénéité et la cohérence de leurs résultats.~»\\
Rapport jury 2024 : «~Le jury attend des réponses détaillées et argumentées. Des résultats donnés directement, sans calcul et sans justification ne répondent pas aux attendus de l'épreuve.~»
\end{tcolorbox}

\end{tcolorbox}

% -----------------------------------------------------------------------
%  PIED - PRESENCE DANS LES SUJETS
% -----------------------------------------------------------------------

\begin{tcolorbox}[
  ficheSection=colPresence,
  title={Présence dans les sujets analysés (2013-2026)},
]
\begin{center}
\footnotesize\renewcommand{\arraystretch}{1.15}
\begin{tabular}{lccccccccccccc}
\toprule
\textbf{Année} & 13 & 14 & 15 & 16 & 17 & 18 & 19 & 20 & 21 & 22 & 23 & 24 & 25 \\
\midrule
E1 (commune)   & $\checkmark$ & $\checkmark$ & $\checkmark$ & - & $\checkmark$ & $\checkmark$ & - & $\checkmark$ & $\checkmark$ & - & - & $\checkmark$ & $\checkmark$ \\
E2 discipl. IE & - & $\checkmark$ & $\checkmark$ & $\checkmark$ & - & - & $\checkmark$ & - & $\checkmark$ & $\checkmark$ & $\checkmark$ & $\checkmark$ & $\checkmark$ \\
\bottomrule
\end{tabular}
\end{center}
\smallskip
\footnotesize\textbf{Sous-thèmes fréquents :}
couple moteur depuis effort de traction (9/13) - rapport de réduction et rendement (8/13) - inertie équivalente ramenée au moteur (5/13) - couple thermique équivalent (4/13) - PFD rotation ($J_{eq}\dot\Omega$) (5/13)

\medskip
\textbf{Ordres de grandeur (repères rapides).}
\begin{itemize}[leftmargin=1.3em]
  \item Accélérations de véhicules en sujet : démarrage $\approx \SI{0,2}{\metre\per\second\squared}$, freinage $\approx \SI{1,5}{\metre\per\second\squared}$ (rame de métro, 2019 E2).
  \item Réducteurs : rendement typique $\eta_{red} \approx 0{,}88$ à $0{,}97$ ; rapport de réduction de $\rho = 0{,}17$ (rame MF67) à $1178{,}8$ (pont levant 2014).
  \item Adhérence rail/roue : $\mu_r \approx 0{,}12$ (2019 E2) ; couple moteur nominal d'une rame MF67 $\approx \SI{885}{\newton\metre}$.
\end{itemize}

\medskip
\textbf{Synthèse express (à refaire sans support).}
\begin{enumerate}[leftmargin=1.4em]
  \item Poser les hypothèses du modèle et les conventions de signe.
  \item Écrire 2 à 4 relations clés avec unités.
  \item Réaliser un calcul numérique complet et vérifier l'ordre de grandeur.
  \item Justifier le résultat par un argument technique de niveau CAPET.
\end{enumerate}

\end{tcolorbox}

\CAPETENDFILE
